\documentclass[12pt,a4paper]{article}

% Required Packages for English and Arabic
\usepackage{fontspec}
\usepackage{polyglossia}
\usepackage{geometry}
\usepackage{titlesec}
\usepackage{xcolor}
\usepackage{enumitem}

% Setup Document Geometry
\geometry{
    a4paper,
    top=1in,
    bottom=1in,
    left=1in,
    right=1in
}

% Setup Languages and Fonts for RTL
\setdefaultlanguage{arabic}
\setotherlanguage{english}
\newfontfamily\arabicfont[Script=Arabic]{Amiri}

% Define colors for headers
\definecolor{primary}{RGB}{30, 64, 175}
\definecolor{secondary}{RGB}{3, 105, 161}

% Formatting Headers
\titleformat{\section}
  {\normalfont\Large\bfseries\color{primary}}{\thesection.}{1em}{}
\titleformat{\subsection}
  {\normalfont\large\bfseries\color{secondary}}{}{0em}{}

% RTL List formatting
\setlist[itemize]{leftmargin=0pt, rightmargin=2em}
\setlist[enumerate]{leftmargin=0pt, rightmargin=2em}

\begin{document}

\begin{titlepage}
    \centering
    \vspace*{3cm}
    
    {\Huge \textbf{\color{primary}نظام مراقبة الزراعة الذكية المعتمد على الذكاء الاصطناعي}\par}
    \vspace{1cm}
    {\Large \textcolor{gray}{\textbf{وثيقة تصميم وتنفيذ النظام}}\par}
    
    \vspace{3cm}
    
    {\Large \textbf{مقدم إلى:}\\ د. محمد الشيشتاوى\par}
    \vspace{2.5cm}
    
    {\Large 13/3/2026\par}
\end{titlepage}

\newpage

\section*{الملخص (Abstract)}
يشهد القطاع الزراعي توجهاً متزايداً نحو تبني التقنيات الحديثة لتحسين الكفاءة، زيادة العائد، وتحقيق الاستدامة. تقدم هذه الوثيقة تصميماً متكاملاً وتنفيذاً لـ "نظام مراقبة الزراعة الذكية المعتمد على الذكاء الاصطناعي". يوفر النظام المُقترح حلاً يعتمد على بيئة الإنتاج الفعلي (production-ready) يدمج البيانات الحية للمستشعرات مع التحليل المخصص لصحة المحاصيل، وتوصيات الري المبنية على الذكاء الاصطناعي. تم تطوير النظام باستخدام تقنية MERN stack، ويدمج تحليلات توقعية باستخدام (الإنحدار الخطي)، ليقدم منصة شاملة وقابلة للتوسع لتلبية احتياجات الزراعة الحديثة.

\vspace{0.5cm}

\section{المقدمة (Introduction)}
تواجه الزراعة الحديثة تحدياً كبيراً في تحسين استخدام الموارد مع الحفاظ على زيادة إنتاجية المحاصيل في ظل ظروف بيئية غير متوقعة. غالباً ما تعتمد طرق الزراعة التقليدية على التقديرات والخبرة الشخصية، مما يؤدي إلى استهلاك مفرط للمياه، تأخر ملحوظ في اكتشاف الأمراض، وتوفير ظروف نمو غير مثالية.

يعالج المنظور المقترح هذه التحديات حصرياً من خلال الاستفادة من إنترنت الأشياء (IoT) والذكاء الاصطناعي لإدخال قرارات دقيقة بالبيانات. صُمم النظام لمراقبة العوامل البيئية مثل الحرارة، الرطوبة، رطوبة التربة، وحموضة التربة، وضبط معايير التحليل بناءً على المحصول المُختار (مثل القمح، الذرة، الطماطم).

\vspace{0.5cm}

\section{أهداف المشروع (Project Objectives)}
الهدف الأساسي من هذا المشروع هو سد الفجوة بين الممارسات الزراعية التقليدية والتطور التكنولوجي. ويهدف المشروع تفصيلياً إلى تطبيق الحلول التالية:
\begin{itemize}
    \item \textbf{تعظيم كفاءة الموارد:} تحسين استخدام المياه والأسمدة لضمان الإمداد الدقيق لاحتياجات المحصول عبر رشاشات ذكية.
    \item \textbf{الاكتشاف الاستباقي للأعطال:} التحول من الزراعة التفاعلية للإدارة الاستباقية لتوقع التغيرات للتنبؤ بموجات الحر أو الجفاف قبل حدوثها.
    \item \textbf{تسهيل التعاون الفعّال:} إنشاء منصة تتيح للمهندس الزراعي إسناد المهام بشكل مباشر ومدروس للمزارعين في الحقل.
    \item \textbf{التخصيص التلقائي:} توفير محرك برمجي يعدل نقاط حساب صحة المحصول ومخاطره بناءً على الفسيولوجيا الخاصة لكل فصيلة نبات.
    \item \textbf{رؤى وبيانات واضحة:} تحويل أرقام الحساسات المعقدة لإحصائيات ونسبة مئوية واضحة (0-100\%).
\end{itemize}

\vspace{0.5cm}

\section{بنية النظام التقنية (System Architecture)}
بُني النظام بالكامل على أحدث معمارية برمجية في الويب (MERN Stack). وهو يعتمد على نموذج العميل - الخادم المزود بمحرك محاكاة داخلي للبيانات.

\subsection{واجهة المستخدم (Frontend)}
\begin{itemize}
    \item \textbf{إطار العمل:} تم استخدام تقنية React (Vite) لتقديم استجابة فائقة السرعة للمستخدم.
    \item \textbf{التصميم (Styling):} يتم الاعتماد على مكتبة TailwindCSS.
    \item \textbf{الرسوم البيانية:} تُستخدم مكتبة Chart.js لتمثيل ورسم اتجاهات البيانات بوضوح.
    \item \textbf{الاتصال اللحظي:} تعمل مكتبة Socket.io-client على جلب القراءات اللحظية دون الحاجة لإعادة تحميل الصفحة.
\end{itemize}

\subsection{الخادم (Backend)}
\begin{itemize}
    \item \textbf{بيئة التشغيل:} يعمل السيرفر بتقنية Node.js مع إطار عمل Express.js.
    \item \textbf{المحرك الحي:} يقوم Socket.io ببث القراءات بشكل دوري كل 5 ثوانٍ لجميع المتصلين بالمنصة.
    \item \textbf{قاعدة البيانات:} تخزن قاعدة MongoDB بيانات الاعتماد، البيانات التاريخية، والإنذارات.
    \item \textbf{المحرك التحليلي:} يتضمن خوارزمية Linear Regression للتوقع المستقبلي المدى.
\end{itemize}

\vspace{0.5cm}

\section{المميزات الأساسية (Key Features)}

\subsection{نظام الصلاحيات (Role-Based Authentication)}
النظام محمي بواسطة JSON Web Tokens وينقسم لـ :
\begin{itemize}
    \item \textbf{المزارع (Farmer):} يتابع الشاشة البسيطة وينهي المهام.
    \item \textbf{المهندس (Analyst/Engineer):} لديه وصول للمؤشرات التحليلية وإسناد الأوامر.
    \item \textbf{المدير (Admin):} يملك كافة الصلاحيات على النظام والواجهة البرمجية.
\end{itemize}

\subsection{المراقبة الحية ومنطق التخصيص للمحصول}
يتدفق دفق البيانات على مدار الساعة. يتم ضبط "تارجت" درجات الحرارة آلياً فور اختيار نوع المحصول من النظام (مثلاً استبدال القمح بالطماطم يغير معطيات الخطر).

\subsection{الري الذكي وتقييمات المستقبل (Predictive Analytics)}
يحلل النظام قراءات التربة لمنح اقتراحات فورية مثل فتح صمامات الري أو الرذاذ فوقي التبريد. الخوارزميات تتنبأ بحرارة الـ 30 دقيقة القادمة مسبقاً.

\subsection{اكتشاف الخلل الآلي (Anomaly Detection)}
تسجل تنبيهات فورية تظهر باللون الأحمر عند خروج أي عامل بيئي عن السيطرة، ما يجبر المشرفين على الالتفات السريع.

\vspace{0.5cm}
\newpage

\section{سيناريو حالة الاستخدام الشامل (Use Case Scenario)}

\subsection{المرحلة الأولى: المراقبة واكتشاف الخلل السريع}
أثناء فصل الصيف، يراقب النظام محصول "الطماطم". فجأة، ترتفع درجة الحرارة لـ 38°C وتنخفض رطوبة التربة بقوة لتصبح 25\%. ولأن النظام مبرمج بأن الطماطم تحتاج لدرجة حرارة تتراوح بين 20-28°C تُسجل هذه التغيرات كـ "حالة شاذة خطيرة".

\subsection{المرحلة الثانية: التنبؤ والإنذار المبكر}
بشكل متزامن، يقوم الذكاء الاصطناعي يتوقع بأن الحرارة ستتجاوز الـ 40°C خلال نصف ساعة. نتيجة لذلك، ينخفض تقييم "صحة المحصول" لـ 65\% تلقائياً ويُطلق النظام إنذاراً حرجاً ومرئياً لجميع لوحات القيادة لحظياً.

\subsection{المرحلة الثالثة: قرار المهندس المعالج}
يلاحظ المهندس الإنذار فوراً. ومن خلال قسم "الري الذكي"، يُقترح عليه تنفيذ عملية (ري متزايد وتبريد محيطي لمدة 45 دقيقة). يُصدر المهندس توصية عاجلة ويُرسلها فوراً في هيئة "مهمة إسناد مباشرة" إلى المزارع.

\subsection{المرحلة الرابعة: التنفيذ الفعلي (المزارع)}
يصل الإشعار إلى المزارع. يعتمد التوصية، يفتح محابس المياه ويبدأ التبريد. بعد انقضاء الوقت المطلوب، تنخفض الحرارة للدرجة المقبولة 32°C وتعود رطوبة التربة للمعدل المطلوب 60\%. يُنهي المزارع المهمة تقنياً في النظام، ليعود مؤشر صحة المحصول للمستوى الأخضر الآمن 88\%.

\subsection{المرحلة الخامسة: إشراف وتصدير البيانات (Admin)}
يدخل مدير النظام ليلاً بحسابه، فيرى الإحصائية لليوم كاملة، فيقوم باستخراج البيانات على هيئة ملف CSV للاحتفاظ بها في الأرشيف ولاستخدامها لتدريب الموديل مستقبلاً بشكل أدق.

\vspace{0.5cm}

\section{الخاتمة والعمل المستقبلي (Conclusion and Future Work)}
أثبت هذا النموذج أن الاعتماد على الذكاء الاصطناعي وإنترنت الأشياء يتجاوز مجرد مراقبة المحاصيل، ليصبح منصة تدير الأزمة قبل وقوعها وتتخذ القرار الصحيح لحماية المزارع. في المستقبل يمكن ترقية النظام عبر:
\begin{itemize}
    \item استخدام الكاميرات لتعرف الذكاء الاصطناعي البصري (Computer Vision via CNNs) على أمراض النباتات من شكل الورقة.
    \item الارتقاء لنماذج تعلم آلي معقدة تعطي تنبؤ موسمي لشهور الطقس بدلاً من مجرد ساعات قليلة.
    \item الربط المباشر مع المعدات (كالجرارات الذكية ومضخات المياه الكبيرة) لتعمل بشكل آلي متكامل بدون أدنى تدخل بشري.
\end{itemize}

\end{document}
